% Cover Letter for PLOS ONE (1-page version)
\documentclass[11pt,letterpaper]{article}
\usepackage[margin=1in]{geometry}
\usepackage{times}
\usepackage{hyperref}
\setlength{\parskip}{6pt}
\setlength{\parindent}{0pt}

\begin{document}

\pagestyle{empty}

\noindent
February 2, 2026

\vspace{6pt}

\noindent
Editor-in-Chief, PLOS ONE

\vspace{6pt}

\noindent
\textbf{Subject:} ``Signal Hierarchical Petri Nets: Formal Semantics of Hierarchical Regulatory Control of Biological Systems''

\vspace{6pt}

Dear Editor,

I submit this manuscript addressing a fundamental barrier in computational biology: no existing formalism can predict cellular commitment thresholds from first principles. Classical models treat metabolites as either substrates (metabolic networks) or switches (gene networks), never capturing molecules functioning simultaneously as metabolic currencies and regulatory signals through modeler designation. For example, when \textit{Bacillus subtilis} enters sporulation, experiments measure commitment at 2.21 $\pm$ 0.18 mM ATP---yet no framework can compute this value without parameter fitting.

\textbf{1. Signal Hierarchy Theory (Original):} We prove two theorems: (1) signal flow graphs must be acyclic (commitment finality), and (2) lower-layer depletion structurally preempts higher-layer transitions (metabolic override).

\textbf{2. Signal Hierarchical Control (Original):} Signal flow arcs implement consumptive commitment semantics distinct from test arcs: token consumption creates irreversible basin boundaries. Layer functions and preemption predicates formalize hierarchical dependencies.

\textbf{3. Signal Hierarchical Petri Nets (New Formalism):} We extend Bio-PN from 5-tuple to 13-tuple with dual arc semantics: normal arcs (mass transfer), signal flow arcs (information propagation). Modelers designate metabolites (ATP, GTP, NADH, cAMP, Ca$^{2+}$) as signals when biological evidence supports regulatory roles.

\textbf{4. Unified Metabolic-Regulatory Modeling (Original):} Thresholds emerge as structural properties: $M_{\text{commit}} = \theta + W_s$. For \textit{B. subtilis} with ATP designated as signal, we predict 2.38 mM (7\% error, zero fitted parameters) using only literature values.

\textbf{5. Validation:} Mathematical proofs establish correctness; biological case demonstrates applicability. This enables synthetic biology circuit design with formal guarantees, metabolite-dependent drug targets, and patient-specific predictions.

Code/data: \url{https://github.com/simao-eugenio/shypn} (MIT, Zenodo). Manuscript: 33 pages with complete proofs and validation.

\vspace{6pt}

\noindent
Sincerely,

\vspace{12pt}

\noindent
\textbf{Eugênio Simão, Ph.D.}\\
Assistant Professor, Computer Science\\
Universidade Federal de Santa Catarina\\
Email: eugenio.simao@ufsc.br

\end{document}
